\documentclass[11pt]{article}
\usepackage[a4paper,margin=1in]{geometry}
\usepackage{amsmath,amssymb,amsthm,mathtools}
\usepackage{booktabs,array,longtable}
\usepackage{enumitem}
\usepackage{microtype}
\usepackage[T1]{fontenc}
\usepackage{lmodern}

\newtheorem{theorem}{Theorem}
\newtheorem{lemma}[theorem]{Lemma}
\newtheorem{proposition}[theorem]{Proposition}
\newtheorem{remark}[theorem]{Remark}
\newcommand{\EE}{\mathsf E}
\newcommand{\V}{\mathbb V}
\newcommand{\F}{\mathcal F}
\newcommand{\nuA}{\nu^{A}}
\newcommand{\nuB}{\nu^{B}}

\title{A strictly positive finite-range continuous-time IPS with a stationary non-quasilocal measure}
\author{}
\date{}

\begin{document}
\maketitle

\begin{abstract}
We give an explicit counterexample on a directed regular tree.  The local state
space is the three-symbol alphabet $\EE=\{0,1,2\}$.  The dynamics is a
continuous-time nearest-neighbour interacting particle system in which an
oriented edge updates its ordered pair of endpoint spins according to a fixed
$9$-state Markov generator.  Every off-diagonal pair transition has a strictly
positive rate, uniformly bounded away from zero.  Nevertheless the dynamics
has two distinct full-support Bernoulli product stationary measures.  Their
nontrivial convex mixture is stationary but has no continuous version of its
one-site conditional probabilities, and hence is not quasilocal.
\end{abstract}

\section{The directed tree and the generator}

Let $\mathcal T=(\V,\vec E)$ be the directed binary tree: every vertex $x$ has
one parent $p(x)$ and exactly two children.  Thus
\[
  \vec E=\{(x,p(x)):x\in\V\}.
\]
The underlying undirected graph is the $3$-regular tree.  Put
$\Omega=\EE^{\V}$ with $\EE=\{0,1,2\}$.

For an oriented edge $e=(x,p(x))$, identify its ordered pair state
$(\eta_x,\eta_{p(x)})\in\EE^2$ with an index in $\{0,\dots,8\}$ via
\[
  \iota(a,b)=3a+b.
\]
For $i\ne j$, let $q_{ij}$ be the rate in Table~\ref{tab:Q}.  A transition
$i\to j$, where $i=\iota(a,b)$ and $j=\iota(c,d)$, replaces the pair
$(\eta_x,\eta_{p(x)})=(a,b)$ by $(c,d)$ and leaves all other spins unchanged.
The formal generator is
\begin{equation}
  (Lf)(\eta)
  =\sum_{x\in\V}\sum_{j\ne \iota(\eta_x,\eta_{p(x)})}
  q_{\iota(\eta_x,\eta_{p(x)}),j}
  \bigl[f(\eta^{x,p(x);j})-f(\eta)\bigr],
  \label{eq:generator}
\end{equation}
for local $f$, where $\eta^{x,p(x);j}$ denotes the configuration obtained by
putting the pair indexed by $j$ on $(x,p(x))$.

\begin{table}[ht]
\centering
\small
\setlength{\tabcolsep}{3.2pt}
\renewcommand{\arraystretch}{1.25}
\caption{Off-diagonal rates $q_{ij}$.  The state order is
$00,01,02,10,11,12,20,21,22$.  A dash denotes the diagonal.}
\label{tab:Q}
\begin{tabular}{c|ccccccccc}
\toprule
 &00&01&02&10&11&12&20&21&22\\
\midrule
00 &--&$\frac{2767709}{9766950000}$&$\frac{10385390113}{33207630000000}$&$\frac{589477}{1953390000}$&$\frac{62846}{203478125}$&$\frac{1813959}{6511300000}$&$\frac{4589347}{19533900000}$&$\frac{2536387}{9766950000}$&$\frac{204841}{1028100000}$\\
01 &$\frac{9391756087}{33207630000000}$&--&$\frac{9595487813}{33207630000000}$&$\frac{4516031}{19533900000}$&$\frac{1545599}{6511300000}$&$\frac{1265009}{6511300000}$&$\frac{5577797}{19533900000}$&$\frac{5102083}{19533900000}$&$\frac{305591}{1028100000}$\\
02 &$\frac{21707821}{74290000000}$&$\frac{13}{50000}$&--&$\frac{13}{50000}$&$\frac{13}{50000}$&$\frac{13}{50000}$&$\frac{13}{50000}$&$\frac{13}{50000}$&$\frac{13}{50000}$\\
10 &$\frac{13}{50000}$&$\frac{630401}{25000}$&$\frac{13}{50000}$&--&$\frac{13}{50000}$&$\frac{13}{50000}$&$\frac{13}{50000}$&$\frac{17029}{20000}$&$\frac{13}{50000}$\\
11 &$\frac{13}{50000}$&$\frac{17809}{20000}$&$\frac{13}{50000}$&$\frac{13}{50000}$&--&$\frac{13}{50000}$&$\frac{13}{50000}$&$\frac{13}{50000}$&$\frac{13}{50000}$\\
12 &$\frac{13}{50000}$&$\frac{13}{50000}$&$\frac{13}{50000}$&$\frac{13}{50000}$&$\frac{13}{50000}$&--&$\frac{13}{50000}$&$\frac{13}{50000}$&$\frac{13}{50000}$\\
20 &$\frac{216521}{50000}$&$\frac{13}{50000}$&$\frac{1411131}{50000}$&$\frac{13}{50000}$&$\frac{13}{50000}$&$\frac{13}{50000}$&--&$\frac{13}{50000}$&$\frac{13}{50000}$\\
21 &$\frac{34733}{100000}$&$\frac{13}{50000}$&$\frac{89301}{100000}$&$\frac{433037}{100000}$&$\frac{13}{50000}$&$\frac{5113}{2000}$&$\frac{13}{50000}$&--&$\frac{13}{50000}$\\
22 &$\frac{13}{50000}$&$\frac{13}{50000}$&$\frac{17039}{100000}$&$\frac{13}{50000}$&$\frac{13}{50000}$&$\frac{733}{100000}$&$\frac{104209}{25000}$&$\frac{13}{50000}$&--\\
\bottomrule
\end{tabular}
\end{table}

Every displayed number is positive.  The smallest one is
\begin{equation}
 \varepsilon_*:=\frac{1265009}{6511300000}>0.
 \label{eq:minrate}
\end{equation}
Hence every possible change of an ordered nearest-neighbour pair is allowed at
rate at least $\varepsilon_*$.  In particular there are no forbidden
single-site changes either: keeping one endpoint fixed while changing the other
is among the off-diagonal pair transitions.

The row sums are finite; the largest is less than $33$.  Every vertex belongs to
exactly three oriented edges.  Consequently the dynamics has finite range and
uniformly bounded local activity.  The standard Harris construction therefore
defines a Feller interacting particle system with generator
\eqref{eq:generator}.  No existence issue for the process will be used below
beyond this elementary bounded-rate construction.

\section{A stationarity criterion on the directed binary tree}

Let $\pi=(\pi_0,\pi_1,\pi_2)$ be a probability vector with every $\pi_a>0$,
and let
\[
  \nu_\pi=\bigotimes_{x\in\V}\pi
\]
be the corresponding product measure.  Put
$w_\pi(a,b)=\pi_a\pi_b$.  For the edge generator $Q=(q_{ij})$, define its
adjoint defect relative to the pair-product law by
\begin{equation}
 h_\pi(a,b)
 :=\frac{1}{\pi_a\pi_b}
   \sum_{(c,d)\ne(a,b)}\pi_c\pi_d\,
     q_{\iota(c,d),\iota(a,b)}
   -\sum_{(c,d)\ne(a,b)}q_{\iota(a,b),\iota(c,d)}.
 \label{eq:defh}
\end{equation}

\begin{lemma}[Tree-divergence criterion]\label{lem:divergence}
Suppose there exists a function $v:\EE\to\mathbb R$ such that
\begin{equation}
  \sum_{a\in\EE}\pi_a v(a)=0
  \label{eq:vmean0}
\end{equation}
and
\begin{equation}
  h_\pi(a,b)=v(b)-2v(a)
  \qquad(a,b\in\EE).
  \label{eq:divergence}
\end{equation}
Then $\nu_\pi$ is stationary for $L$.
\end{lemma}

\begin{proof}
Fix a local function $f$, and let $A\Subset\V$ contain its support.  For an
oriented edge $e=(x,p(x))$, let $K_e$ denote the corresponding summand of the
generator.  A finite-state change-of-variables computation gives
\begin{equation}
 \int K_e f\,d\nu_\pi
 =\int f(\eta)\,h_\pi(\eta_x,\eta_{p(x)})\,d\nu_\pi(\eta).
 \label{eq:adjointidentity}
\end{equation}
Indeed, for a fixed target pair $(a,b)$, the incoming contribution is
$\sum_{(c,d)\ne(a,b)}\pi_c\pi_d q_{(c,d),(a,b)}$, while division by the
pair weight $\pi_a\pi_b$ gives the first term in \eqref{eq:defh}; the outgoing
part gives the second term.

If neither endpoint of $e$ lies in $A$, then $K_e f=0$.  Thus only finitely many
edges occur.  Insert \eqref{eq:divergence} into
\eqref{eq:adjointidentity}.  If $x\in A$ but $p(x)\notin A$, integration over
the outside parent and \eqref{eq:vmean0} reduce the factor
$v(\eta_{p(x)})-2v(\eta_x)$ to $-2v(\eta_x)$.  If
$x\notin A$ but $p(x)\in A$, integration over the outside child reduces it to
$v(\eta_{p(x)})$.  For an internal edge both terms remain.

We may therefore collect the resulting terms by vertices of $A$.  Every
$x\in A$ occurs exactly once as the child endpoint of its parent edge, producing
$-2v(\eta_x)$, and exactly twice as the parent endpoint of its two child edges,
producing $+v(\eta_x)+v(\eta_x)$.  Hence all terms cancel identically:
\[
 \int Lf\,d\nu_\pi
 =\int f\sum_{x\in A}\bigl[-2v(\eta_x)+2v(\eta_x)\bigr]d\nu_\pi=0.
\]
Since this holds for every local $f$, $\nu_\pi$ is stationary.
\end{proof}

\section{Two distinct product stationary measures}

Define
\begin{equation}
  \pi^A=\left(\frac{23}{100},\frac{14}{100},\frac{63}{100}\right),
  \qquad
  \pi^B=\left(\frac{19}{100},\frac{31}{100},\frac12\right),
  \label{eq:pis}
\end{equation}
and the two product measures
\[
  \nuA:=\nu_{\pi^A},\qquad \nuB:=\nu_{\pi^B}.
\]
For $\pi^A$ set
\begin{equation}
 v^A=
 \left(
 -\frac{45787}{3680},
 \frac{87917}{100000},
 \frac{434699}{100000}
 \right),
 \label{eq:vA}
\end{equation}
and for $\pi^B$ set
\begin{equation}
 v^B=
 \left(
 -\frac{12245721}{950000},
 \frac{44491}{50000},
 \frac{21733}{5000}
 \right).
 \label{eq:vB}
\end{equation}
Direct exact arithmetic gives
\begin{equation}
 \sum_a\pi^A_a v^A(a)=0,
 \qquad
 \sum_a\pi^B_a v^B(a)=0.
 \label{eq:meanszero}
\end{equation}
For completeness, the two identities expand to
\[
 \frac{23}{100}\left(-\frac{45787}{3680}\right)
 +\frac{14}{100}\frac{87917}{100000}
 +\frac{63}{100}\frac{434699}{100000}=0,
\]
and
\[
 \frac{19}{100}\left(-\frac{12245721}{950000}\right)
 +\frac{31}{100}\frac{44491}{50000}
 +\frac12\frac{21733}{5000}=0.
\]

The rate table was chosen so that, with $h_\pi$ defined by
\eqref{eq:defh},
\begin{equation}
 h_{\pi^A}(a,b)=v^A(b)-2v^A(a),
 \qquad
 h_{\pi^B}(a,b)=v^B(b)-2v^B(a)
 \label{eq:two-balance}
\end{equation}
for all nine pairs $(a,b)\in\EE^2$.  There is no limiting or approximate
statement here: these are identities of rational numbers.  To make the finite
arithmetic completely inspectable, Tables~\ref{tab:hA} and \ref{tab:hB} list
the values obtained from the left-hand side of \eqref{eq:defh}; they agree
entrywise with the right-hand side of \eqref{eq:two-balance}.

\begin{table}[ht]
\centering
\caption{Exact values of $h_{\pi^A}(a,b)$ computed from Table~\ref{tab:Q}.}
\label{tab:hA}
\begin{tabular}{c|ccc}
\toprule
$a\backslash b$&0&1&2\\
\midrule
0&$\frac{45787}{3680}$&$\frac{59255841}{2300000}$&$\frac{67231827}{2300000}$\\
1&$-\frac{32661057}{2300000}$&$-\frac{87917}{100000}$&$\frac{51773}{20000}$\\
2&$-\frac{48613029}{2300000}$&$-\frac{781481}{100000}$&$-\frac{434699}{100000}$\\
\bottomrule
\end{tabular}
\end{table}

\begin{table}[ht]
\centering
\caption{Exact values of $h_{\pi^B}(a,b)$ computed from Table~\ref{tab:Q}.}
\label{tab:hB}
\begin{tabular}{c|ccc}
\toprule
$a\backslash b$&0&1&2\\
\midrule
0&$\frac{12245721}{950000}$&$\frac{25336771}{950000}$&$\frac{3577589}{118750}$\\
1&$-\frac{13936379}{950000}$&$-\frac{44491}{50000}$&$\frac{32087}{12500}$\\
2&$-\frac{20504261}{950000}$&$-\frac{390169}{50000}$&$-\frac{21733}{5000}$\\
\bottomrule
\end{tabular}
\end{table}

For example, the $(a,b)=(0,0)$ entry in the first table equals
$v^A(0)-2v^A(0)=-v^A(0)=45787/3680$, and the same elementary check applies to
all remaining entries.  Thus \eqref{eq:meanszero} and
\eqref{eq:two-balance} verify the hypotheses of Lemma~\ref{lem:divergence}
for both product laws.

\begin{proposition}\label{prop:twostationary}
Both $\nuA$ and $\nuB$ are stationary for the same generator $L$.
\end{proposition}

\begin{proof}
Apply Lemma~\ref{lem:divergence} first with $(\pi,v)=(\pi^A,v^A)$ and then
with $(\pi,v)=(\pi^B,v^B)$.
\end{proof}

\section{The stationary mixture is not quasilocal}

Set
\begin{equation}
  \mu:=\frac12\nuA+\frac12\nuB.
  \label{eq:mixture}
\end{equation}
By Proposition~\ref{prop:twostationary} and linearity, $\mu$ is stationary.
We now prove that it is not quasilocal.

Fix a vertex $o\in\V$ and enumerate $\V\setminus\{o\}$ as
$x_1,x_2,\dots$.  Define the exterior tail event
\begin{equation}
 T_A:=\left\{\eta:
 \lim_{n\to\infty}\frac1n\sum_{k=1}^n
 \mathbf 1_{\{\eta_{x_k}=0\}}=\frac{23}{100}
 \right\}.
 \label{eq:TA}
\end{equation}
The strong law of large numbers gives
\begin{equation}
 \nuA(T_A)=1,
 \qquad
 \nuB(T_A)=0,
 \label{eq:tailseparates}
\end{equation}
because the corresponding $0$-densities are $23/100$ and $19/100$.
Since $T_A\in\F_{\V\setminus\{o\}}$, the exterior configuration determines
which component of the mixture is present, up to a $\mu$-null set.

\begin{lemma}\label{lem:conditional}
A version of the one-site conditional probability of the event $\{\eta_o=0\}$
is
\begin{equation}
 G(\eta)
 =\frac{23}{100}\,\mathbf 1_{T_A}(\eta)
 +\frac{19}{100}\,\mathbf 1_{T_A^c}(\eta).
 \label{eq:conditional}
\end{equation}
In particular,
\[
 G=\frac{23}{100}\quad\nuA\text{-a.s.},
 \qquad
 G=\frac{19}{100}\quad\nuB\text{-a.s.}
\]
\end{lemma}

\begin{proof}
Let $C\in\F_{\V\setminus\{o\}}$.  Since each component is a product measure,
$\eta_o$ is independent of $C$ under that component.  Hence
\[
 \int_C\mathbf 1_{\{\eta_o=0\}}\,d\mu
 =\frac12\frac{23}{100}\nuA(C)
  +\frac12\frac{19}{100}\nuB(C).
\]
Using \eqref{eq:tailseparates}, the right-hand side is exactly
$\int_C G\,d\mu$.  This is the defining property of the conditional
expectation.
\end{proof}

\begin{theorem}\label{thm:nonquasi}
The stationary measure $\mu$ in \eqref{eq:mixture} is not quasilocal.  More
precisely, the conditional probability
$\mu(\eta_o=0\mid\F_{\V\setminus\{o\}})$ has no continuous version in the
product topology.
\end{theorem}

\begin{proof}
Assume for contradiction that a continuous function
$g:\EE^{\V\setminus\{o\}}\to[0,1]$ is a version of the conditional
probability.  By Lemma~\ref{lem:conditional},
\begin{equation}
 g=\frac{23}{100}\quad\nuA\text{-a.s.},
 \qquad
 g=\frac{19}{100}\quad\nuB\text{-a.s.}
 \label{eq:two-values}
\end{equation}
(where the product measures are understood on the exterior coordinates).

Both $\nuA$ and $\nuB$ have full topological support because every coordinate
probability in \eqref{eq:pis} is strictly positive.  Therefore any set of full
$\nuA$-measure is dense, and any set of full $\nuB$-measure is dense: otherwise
its complement would contain a nonempty cylinder, but every nonempty cylinder
has positive product probability.

Thus the set on which $g=23/100$ is dense, and the set on which $g=19/100$ is
dense.  Continuity forces $g\equiv23/100$ from the first dense set and
$g\equiv19/100$ from the second.  Since $23/100\ne19/100$, this is impossible.
Hence no continuous version exists.  Failure of continuity for one one-site
conditional probability already implies failure of quasilocality.
\end{proof}

\section{Conclusion and scope}

Combining the preceding results gives the promised counterexample.

\begin{theorem}[Counterexample]\label{thm:main}
There exists a continuous-time interacting particle system with
\begin{enumerate}[label=(\roman*)]
\item finite local state space $\EE=\{0,1,2\}$;
\item nearest-neighbour, hence finite-range, transition rules;
\item every off-diagonal local pair transition allowed at a rate at least
$\varepsilon_*=1265009/6511300000$;
\item uniformly bounded local total rates;
\end{enumerate}
which admits a stationary probability measure $\mu$ that is not quasilocal.
\end{theorem}

\begin{proof}
Items (i)--(iv) follow from the construction and Table~\ref{tab:Q}.
Stationarity of $\mu$ follows from Proposition~\ref{prop:twostationary}, and
non-quasilocality follows from Theorem~\ref{thm:nonquasi}.
\end{proof}

\begin{remark}[Single-site versus block-update dynamics]
The counterexample uses nearest-neighbour \emph{pair updates}.  This is a
standard finite-range IPS generator, and it satisfies a stronger local
irreducibility statement than merely allowing every single-site spin change:
every off-diagonal state change of the entire ordered edge pair is allowed with
a uniform positive lower bound.  If one imposes the additional restriction that
\emph{only one site may change at each jump}, the construction above does not
by itself settle that narrower subclass.
\end{remark}

\section*{Exact-arithmetic audit}

For ease of independent checking, the proof uses only the following finite
identities in addition to elementary probability:
\begin{enumerate}
\item all $72$ off-diagonal entries in Table~\ref{tab:Q} are positive, with
minimum $1265009/6511300000$;
\item the two weighted means in \eqref{eq:meanszero} vanish;
\item substituting Table~\ref{tab:Q} into \eqref{eq:defh} yields exactly the
nine rational numbers in Table~\ref{tab:hA} for $\pi^A$ and exactly the nine
rational numbers in Table~\ref{tab:hB} for $\pi^B$;
\item each entry of Table~\ref{tab:hA} equals $v^A(b)-2v^A(a)$, and each entry
of Table~\ref{tab:hB} equals $v^B(b)-2v^B(a)$.
\end{enumerate}
These are finite rational equalities; no approximation or computer-assisted
inequality is part of the argument.

\end{document}
